% Copyright (c) 2026 by C. Klukas.
% Licensed under the MIT License. See LICENSE for details.

\subsection{Mobile agents, again}\label{sec:introduction-agents}

A mobile agent is a program that moves itself, with its state, to the
machine where the data or the resource is, works there, and moves on.
The case for it was made in the 1990s: fewer round trips over slow
links, work that continues while the client is disconnected, services
that are extended by code that arrives rather than by code that is
installed \cite{lange1999seven,harrison1995mobileAgents}. Aglets, the
system that fixed the vocabulary, made agents Java objects with a
lifecycle of creation, cloning, dispatch, deactivation and disposal,
each with a callback, reached through proxies that forward messages to
the agent wherever it is \cite{lange1997aglets,lange1998aglets}.

The first generation had three problems that were design problems, not
accidents. An agent moved as a serialized Java object, so what travelled
was decided by what the author marked serializable, and the classes had
to be present on the destination or fetched from a code base. The Java
virtual machine does not expose the state of a thread, so execution could
not move: the agent's \code{run()} started again from its beginning on
the new host, the weak mobility of Fuggetta et al.\
\cite{fuggetta1998mobility}. And a host ran foreign bytecode inside its
own virtual machine, protected by the Java security model alone, which
was the doubt raised by Harrison, Chess and Kershenbaum before the
systems existed \cite{harrison1995mobileAgents}; the design space of
code shipping, state capture and security is surveyed by Karnik and
Tripathi \cite{karnik1998design}. Serialization that captures too much
or too little, no execution state and no sandbox that one could trust
with a stranger's code: these were the reasons to prefer remote calls,
and remote calls won.

\subsection{What WebAssembly changes}\label{sec:introduction-wasm}

WebAssembly \cite{haas2017wasm} changes the trade-offs of each of the
three. It is a portable binary instruction set with a formal semantics,
compiled from C and C++ by clang and from other languages by their own
toolchains. An instance of a module has one linear memory, a contiguous
array of bytes that only the instance's own code and its host can read
or write, and a set of globals; a C++ program compiled for it keeps its
heap and its shadow stack in that memory. Between two calls into the
instance, the memory and the mutable globals are its complete state, and
the specification defines both without reference to the host, the CPU
or the operating system. The instance is sandboxed by construction: it
reaches the outside only through the functions its host provides as
imports, and a module that imports anything else cannot be
instantiated. Small
runtimes embed in a C program; \paglets{} uses WAMR \cite{wamr}. What
WebAssembly does not give is execution state across a call: a running
function's stack is the runtime's, not the module's. The design of
\paglets{} accepts this and moves agents only between calls.

\subsection{The idea}\label{sec:introduction-idea}

A \paglet{} is a WebAssembly module, named by the SHA-256 hash of its
bytes, together with one instance of it. Its whole state is the
instance's linear memory and exported mutable globals; a \emph{memory
image} of these, the non-zero 64\,KB pages with their hashes and the
globals, is all a host needs to store it, copy it or send it to another
host. A \paglet{} is \emph{created} from a module, with start arguments
and the capabilities its creator passes. It \emph{handles} messages,
requests and timers, one call at a time, through handlers that must
return within a time budget. It is \emph{dispatched} to another host by
a transfer ticket that names a host, a label, a service offer or
\code{any}; the destination fetches the module by hash if it lacks it
and receives the pages it does not hold, and the \paglet{} continues
there with its next message. It is \emph{cloned}, on its host or onto
another one, with a copy of its memory. It is \emph{deactivated} to an
image and \emph{activated} by the next message, timer or wake-up, and
on a host with a state directory it resumes from its last image after
the host restarts. It is
\emph{disposed}, or it ends by a trap. Every transition is a hook the
\paglet{} may implement, and every failure of a transition is a report
it receives as a message. The host that runs it, \code{paglets-host},
keeps the instance in a worker process that the operating system
confines, and belongs to a \emph{mesh}: hosts that share a signed
ledger of members, trusted modules and policy, and nothing else.

\subsection{Design principles}\label{sec:introduction-principles}

Five principles shape the system.

\begin{itemize}
\item \emph{State is the memory image.} No \paglet{} serializes
  anything: a \code{std::vector} in a member moves as the bytes it
  occupies, and nothing in memory can be left behind, because there is
  no list of what to take. The price is a fixed memory limit
  (16\,MB) and images that keep every page that was ever non-zero
  (Sections~\ref{sec:movement} and \ref{sec:discussion}).
\item \emph{Events between handlers.} The host calls into an instance
  from one thread at a time and takes images only between calls. A move,
  a clone or a deactivation requested inside a handler takes effect
  after the handler returns, in a fixed order, and a hook is a call like
  any other. Failures arrive as reports (\code{on\_move\_failed},
  \code{on\_undelivered}, \code{on\_recovered}), not as exceptions
  (Section~\ref{sec:lifecycle}).
\item \emph{Capabilities, not addresses.} A \paglet{} has no address.
  It reaches another \paglet{}, a service or a resource only through a
  capability in a table the host keeps for it: an endpoint, a reply
  capability, a timer, a directory or an artifact. Capabilities move
  with the \paglet{} under a rule per kind, and location records let an
  endpoint find its target after a move (Sections~\ref{sec:model} and
  \ref{sec:communication}).
\item \emph{A signed mesh.} Admins, hosts, owners and signers hold
  Ed25519 keys. Membership, module trust, policy and grants are signed
  records in a grow-only ledger spread by gossip; hosts talk inside
  mutually authenticated, end-to-end encrypted channels; and a
  \paglet{} reaches host resources only through system \paglet{}s, within
  the policy (Section~\ref{sec:mesh}).
\item \emph{Verification in the same code base.} The fault-tolerance
  protocol of zombie \paglet{}s is a pure C++ transition function. The
  guest SDK runs it, an explicit-state model checker in the repository
  explores it, a simulator runs it under time and random faults, and
  live traces are checked against its transition table
  (Section~\ref{sec:zombies}).
\end{itemize}

\subsection{Contributions}\label{sec:introduction-contributions}

\begin{enumerate}
\item A programming model and an ABI (v1, minor version 1.2,
  Appendix~\ref{app:abi}) for mobile agents as WebAssembly instances:
  fourteen imports, six required exports, seven event kinds and three
  reports, covered by 41 conformance tests.
\item A host runtime with sandboxed worker processes on Linux and macOS
  (Windows experimental), whose memory images are independent of the
  CPU and the operating system: one build of the counter module (an
  ABI~v0 build of the feasibility milestone) was moved in four transfers between macOS arm64, Linux arm64 and an x86-64
  build under Rosetta~2 (Section~\ref{sec:evaluation}).
\item A move protocol between hosts that sends only the pages the
  destination lacks (2.2 pages on average on repeated moves of a 4\,MB
  \paglet{}, after 66 on the first), with location records on three
  hosts chosen by consistent hashing, so that endpoints keep working
  across moves (Sections~\ref{sec:movement} and \ref{sec:communication}).
\item A mesh bound only by a signed ledger: four key roles, quorums,
  gossip, Noise channels inside HTTPS, allow/ask/deny policy with
  grants, system \paglet{}s as the only path to resources
  (Section~\ref{sec:mesh}), and relays for hosts without an inbound port
  (Section~\ref{sec:communication-relays}).
\item Zombie \paglet{}s, whose protocol is checked in the repository:
  four invariants and two recoverability goals hold in ten two-replica
  worlds explored completely (up to 22.4 million states), the invariants
  hold to depth 11 in two three-replica worlds, 10\,000 timed runs with
  119\,289 faults show no failure, and a live run with a failure
  schedule on three machines conforms to the specification
  (Section~\ref{sec:zombies}).
\item Tools that stay true to the sources: interface descriptions inside
  the module, a command line and a web UI built from them, and a
  generator that renders the code excerpts and result tables of the
  documentation and of this paper from the sources, checked by a test
  (Section~\ref{sec:tools}).
\item An evaluation on three real machines and in a test suite of 252
  cases run in two modes: request round trips of 1.16--5.16\,ms mean,
  file transfers of 20--40\,MB/s, 100\,000 hexadecimal digits of $\pi$
  in 313\,s on three hosts, and what is not measured
  (Section~\ref{sec:evaluation}).
\end{enumerate}

\subsection{The paper}\label{sec:introduction-outline}

Section~\ref{sec:background} places \paglets{} among Aglets, weak and
strong mobility, process migration, actors and earlier Wasm snapshots,
and names the Python edition of \paglets{}, whose application note
\cite{klukas2026pagletsNote} describes a system that moves a declared
state record between hosts sharing the code; the C++ edition
\cite{pagletsCpp,pagletsCppDocs} keeps its concepts and none of its
formats. Section~\ref{sec:model} defines the model: module and instance,
identity, messages, typed payloads, capabilities and the import set.
Section~\ref{sec:lifecycle} gives the states, the hooks and the order in
which the host runs them. Section~\ref{sec:movement} describes the
memory image, what does not move, transfer tickets and the two-host
protocol. Section~\ref{sec:communication} shows how an endpoint reaches
a \paglet{} that has moved. Section~\ref{sec:mesh} describes the ledger,
passports, policy and system \paglet{}s. Section~\ref{sec:zombies}
presents zombie \paglet{}s and what was verified about them.
Section~\ref{sec:tools} covers interface descriptions, the command line,
the web UI and the documentation generator. Section~\ref{sec:evaluation}
reports the measurements, Section~\ref{sec:discussion} the costs and
limits, and Section~\ref{sec:conclusion} concludes. The appendices list
the ABI (\ref{app:abi}), the complete example \paglet{}s
(\ref{app:examples}), the hooks and ordering guarantees
(\ref{app:events}) and how to reproduce the results
(\ref{app:reproduce}).
